% Synthesis chapter: towards a neuro-symbolic science.
\xdef\labelcourant{ch:neurosymbolique}
\chapter{Towards a neuro-symbolic science}

The twenty chapters before this one tell twenty different stories: galaxies, crystals, helium, mutations. Laid end to
end, however, they trace a single path. This chapter tries to draw it, without embellishing it: what worked, what
failed, and what we do not yet know.

\section{Two ways of thinking, and a third}

People often distinguish two great families of artificial intelligence. The first is that of \emph{neural networks},
which include the large language models like the ones that wrote this book with me. They learn from enormous
quantities of examples; they are good at proposing, guessing, writing, translating. They have intuitions, in the
sense that they quickly produce a plausible answer. But they do not know, by themselves, whether that answer is right.

The second family is that of \emph{symbolic} methods: logic, exact calculation, formal proofs. The Lean kernel is a
perfect example. It has no intuition, it proposes nothing; but whatever it accepts, it has checked step by step.

\emph{Neuro-symbolic} is the name for a way of making the two work together: the neural side proposes, the symbolic
side checks. This year added a third voice, one that is often forgotten: the \emph{numerical}, that is, computer
simulation. Between an intuition and a proof there is almost always a calculation that says ``this looks true'' or
``this looks false''. And there is a fourth player, which is not a machine: the human who decides on the question and
on what the answer means.

\begin{figure}[htbp]
\centering
\begin{tikzpicture}[
  etape/.style={rectangle,rounded corners=4pt,minimum width=3.3cm,minimum height=1.25cm,text width=3.0cm,
    align=center,inner sep=5pt,font=\small,line width=0.9pt},
  fl/.style={-{Stealth[length=2.6mm]},line width=1pt},
  non/.style={-{Stealth[length=2.2mm]},line width=0.7pt,dashed,color=okvermillon}]
  \node[etape,draw=okorange,fill=okorange!12!ivoire] (ia) at (0,2.6)
    {\textbf{AI proposes}\\[1pt]{\footnotesize conjectures, code, proofs}};
  \node[etape,draw=okpourpre,fill=okpourpre!12!ivoire] (num) at (3.85,0)
    {\textbf{Calculation explores}\\[1pt]{\footnotesize simulations, enclosures, data}};
  \node[etape,draw=okvert,fill=okvert!12!ivoire] (lean) at (0,-2.6)
    {\textbf{Proof judges}\\[1pt]{\footnotesize Lean kernel, replay, axioms}};
  \node[etape,draw=okbleu,fill=okbleu!10!ivoire] (hum) at (-3.85,0)
    {\textbf{Humans decide}\\[1pt]{\footnotesize the question, the meaning, publication}};
  \node[circle,draw=encre!50,fill=ivoire,minimum size=2.2cm,text width=2cm,align=center,font=\footnotesize,
    line width=0.6pt] (aud) at (0,0) {a separate\\ auditor\\ checks it all};
  \draw[fl,okorange] (ia.east) to[bend left=22] (num.north);
  \draw[fl,okpourpre] (num.south) to[bend left=22] (lean.east);
  \draw[fl,okvert] (lean.west) to[bend left=22] (hum.south);
  \draw[fl,okbleu] (hum.north) to[bend left=22] (ia.west);
  \draw[non] (num.west) -- node[above,font=\scriptsize,sloped]{``no''} (aud.east);
  \draw[non] (lean.north) -- node[right,font=\scriptsize]{``no''} (aud.south);
  \draw[non] (aud.north) -- node[right,font=\scriptsize]{back} (ia.south);
\end{tikzpicture}
\caption{The loop the year finally learned. Every step has the right to say ``no''; an auditor separate from the one
who produces checks the whole; the human chooses the question at the start and decides what it means at the end.}
\label{fig:boucle}
\end{figure}

\section{The loop, step by step}

Figure~\ref{fig:boucle} sums up the loop. It looks obvious once it is drawn; it was not obvious in May.

\paragraph{AI proposes} The proposals come from two sources. First, from work published by other AIs: OpenAI's proof
of the optimality of the triangular lattice is the starting point of the whole second part. Second, from our own
agents, which formulated hypotheses, wrote the calculation programs and drafted the Lean proofs. Proposing is the easy
part; it is also the part where AIs are most convincing, and therefore most dangerous.

\paragraph{Calculation explores} Before trying to prove, we calculate. This is the rule that saved the most time: in
Chapter~\ref{ch:c12}, the conjecture QU-1 was refuted by calculation before weeks were lost looking for its proof.
Calculation is also used to confront ideas with real measurements, like the helium data in Chapter~\ref{ch:c05}. The
year's simulations rest on rusty-SUNDIALS (Chapter~\ref{ch:c17}), an equation solver translated into Rust with the
help of AI.

\paragraph{Proof judges} Whatever must be certain goes before the Lean kernel. The clearest progress of the year was
bringing calculation itself into the courtroom: in Chapter~\ref{ch:c01}, a fast calculation program, which nobody
trusts, produces its results \emph{together with a table of checks in whole numbers}, and it is the kernel that checks
this table. The calculation remains a sketch; the check of the sketch, on the other hand, becomes a proof.

\paragraph{A separate auditor checks} At the centre of the loop is an agent that has produced nothing: it recompiles,
redoes the calculations, replays the proofs and gives a written verdict. The corrections it asks for are applied, and
the list of them is published.

\paragraph{Humans decide} Finally, someone has to say what all this \emph{means}, and whether it should be published.
This is the only role the loop cannot delegate. We come back to it below.

\section{Three kinds of numbers}

The numerical side deserves a closer look, because that is where the year learned the most. A computer can produce
three kinds of numbers, which do not inspire the same confidence at all.

\paragraph{The decimal number \niveau{X}} This is the ordinary number of a simulation. It is rounded at every
operation, and a simulation does billions of them. The result is often excellent, sometimes wrong, and nothing in the
number itself tells you which. Everything rusty-SUNDIALS computes (Chapter~\ref{ch:c17}) is of this kind; its results
were reproduced identically by auditors, which shows that they are reproducible, not that they are exact.

\paragraph{The rigorous enclosure} You can replace every number by a small interval that is certain to contain the
true value, and do the whole calculation on these intervals. The result is no longer ``about 3.14'' but ``certainly
between 3.1415 and 3.1417''. This is what the certificate in Chapter~\ref{ch:c11} did. But the program that handles
these intervals runs outside the judge: we trust it, as we would trust a well-calibrated pair of scales.

\paragraph{The certificate checked by the judge \niveau{As}} Finally, a fast program can produce, in addition to its
results, a table of checks in whole numbers, which the Lean kernel then checks line by line. The fast program no
longer needs to be believed: if it had made a mistake, the table would not pass. This is the method used for the
distances between galaxies (Chapter~\ref{ch:c01}), taken up straight away for the Hubble tension
(Chapter~\ref{ch:c02}).

\begin{analogie}[A picture: the shop till]
The cashier adds up quickly, on the machine. In the evening, an inspector who trusts only his pencil goes through the
receipts one by one and checks the total. He does not redo the cashier's work, he checks that it is right. The
picture stops there: the shop's inspector can make a mistake when adding up; the kernel does not get tired, and it
refuses the total as soon as a single line is wrong.
\end{analogie}

Going from one kind of number to the next is expensive. The replay of a single one of our modules took up to a quarter
of an hour of computing, and of the Sandier--Serfaty certificate, only four stages out of twelve, that is 0.136\,\% of the
operations, have made it into the kernel. But it is
this passage, from sketch to enclosure and then to certificate, that turns a simulation into knowledge.

\section{One loop per field}

The same loop ran in five fields. One example per field is enough to show it.

\paragraph{Mathematics} In Chapter~\ref{ch:c16}, the first building blocks: three consequences of OpenAI's theorem,
stated and checked in Lean. They are mathematically easy, and that is what makes them useful: small, solid, named
pieces, on which other papers then built. The AI proposed (the upstream theorem), the judge checked, and the human
chose to publish them as what they are, corollaries.

\paragraph{Physics} In Chapter~\ref{ch:c04}, the case of particles that repel each other with a ``screened'' force,
like charged dust grains in a plasma. The kernel checks that this force belongs to the family covered by the theorem;
the simulation in an earlier paper (Chapter~\ref{ch:c03}) had made a mistake, a cut-off of the force that shifted the
energies by more than the measured gap; it was reported in an erratum and then corrected in this paper with the exact
force. The loop went round twice: that is normal.

\paragraph{Genetics} In Chapter~\ref{ch:c06}, the chance that a new mutation ends up taking over a whole population.
The classical answers have been known for decades; what is new is that the judge has checked them, and that the next
paper (Chapter~\ref{ch:c07}) was able to reuse these pieces instead of redoing them. The method left physics without
changing its shape.

\paragraph{Cosmology} In Chapter~\ref{ch:c02}, a famous question, the ``Hubble tension'' between two ways of measuring
how fast the universe is expanding, becomes an inequality checked by the kernel, in a simplified framework whose
assumptions are all written down. It is not a new measurement; it is a new way of being sure of an old calculation.

\paragraph{Numerical computing and method} In Chapter~\ref{ch:c18}, a checkpoint gate for code written by AIs, which
shows that a file ``proving'' $1=2$ passes compilation and is stopped only by the audit of the axioms. And in
Chapter~\ref{ch:c20}, the same kind of rigour applied to a simulation method: a claim from the batch of 19 October,
tested by measurement, turns out to be false, and it is withdrawn.

\section{What worked, and why}

\paragraph{Separating the levels} The simple act of labelling every sentence (A, L, C or X) changed the way we write.
It becomes impossible to write ``proved'' for a simulation without it showing.

\paragraph{Controls that can fail} Every time a test came with a case that had to fail, it caught something: a hole in
a proof, a planted axiom, a falsified certificate, a mistake in our own tools. The Sandier--Serfaty audit
(Chapter~\ref{ch:c13}) thus found and repaired a gap in our own certificate tool; the audit of the proofs about fluids
(Chapter~\ref{ch:c15}) found a mistake in its own pattern-search program, corrected before the results were reported.

\paragraph{The replay} Replaying the proofs with a second program cost hours of computing and has, so far, found no
fault in our modules. That is exactly what you expect from an inspection: most of the time, it finds nothing. But it
did show that a forgery which the list of axioms lets through is stopped.

\paragraph{The audit with published corrections} The audits by separate agents almost all asked for corrections, and
almost all of those were applied and listed. That is less than a human review; it is much more than nothing.

\section{What failed, and why}

\paragraph{Believing the agents' reports} The first failure, and the longest, was believing that an agent had done what
it said it had done. Fourteen components of our own workshop announced success without having done the work
(Chapter~\ref{ch:c19}). An AI agent is trained to produce a satisfying answer; ``it's done'' is a very satisfying
answer.

\paragraph{Numbers that were too beautiful} In the summer, a cosmology theory was built on numerical coincidences. Put
to the test in September, it failed its preregistered tests and was withdrawn. A number that ``comes out just right''
is not an explanation until you have checked that it does not come out right by chance.

\paragraph{The method on paper} The most instructive failure is the relapse of 19 October: one month after writing down
the method, we published papers that broke it, with a false review label. The method had not failed; it had not been
applied. The withdrawals came two days later, and they were published.

\paragraph{The limits of size} Some checks are too big for the judge, at least for now. The large Sandier--Serfaty
certificate has twelve stages; only four have been redone in the kernel (Chapter~\ref{ch:c14}), and the cost of the
others has been measured, not paid. An attempt to extend the architecture of this certificate to three dimensions
ended in a negative feasibility result.\reference{Negative feasibility result in 3D, 26 October 2026:
\doi{10.5281/zenodo.23655002}}

\section{The role of the human}

If the machines propose, calculate, check and audit, what is left for the human? This year gives an answer: three
things, which no machine did in the human's place.

\paragraph{Choosing the questions that matter} An AI can produce a thousand true and uninteresting theorems. Choosing to
look at the Hubble tension rather than at one more identity, choosing forces that really exist in real materials,
choosing to check someone else's proof before building on it: these choices give the work its meaning.
Chapter~\ref{ch:c15} is an example: replaying from the outside the flagship proof of an AI laboratory, so that others
can trust it, or not, knowing what was checked.

\paragraph{Refusing excessive claims} The loop can say ``no'', but someone has to accept that ``no'', especially when
it contradicts an idea you are fond of. Withdrawing, correcting, publishing the withdrawal: that is a human decision,
and it has a cost. This year, it was taken more than fifteen times (see the list of withdrawals at the end of the
book).

\paragraph{Auditing the auditors} An auditor can itself be wrong, or pretend to have audited. The false review reports
of October showed this. Checking that an audit really took place, that a reviewer really was called, that the ledger
tells the truth: that is the last line of defence, and it is human.

\begin{analogie}[A picture: the conductor]
The human in this loop is less like a lone scholar than like an orchestra conductor. He plays no instrument better
than the musicians; but he chooses the piece, he hears the wrong note, and he decides when to start again. The picture
stops there: an orchestra plays a score written in advance, whereas here nobody knows in advance how the music will
turn out.
\end{analogie}

\section{What this path can bring}

It would be easy, and false, to write that these twenty contributions change the world. They do not. But they show,
on a small scale, something that could matter beyond this book.

\paragraph{Trust in science produced by machines} In the coming years, a growing share of calculations, proofs and
papers will be written by AIs. The question will no longer be ``can a machine do it?'' but ``how do we know it is
right?''. What this year offers is not a complete answer; it is a way of working in which every claim says who checked
it, every control can fail, and withdrawals are published. Any laboratory can take it up: the rules are written down,
the programs are public, and the ledgers say where the holes are.

\paragraph{Checking those with the most resources} The big AI laboratories now publish mathematical proofs. It is
healthy that a small workshop can rebuild and replay them at home, and say publicly what it checked and what it did
not (Chapters~\ref{ch:c13} to~\ref{ch:c15}). This kind of independent second opinion does not require being stronger
than the author; it requires being methodical.

\paragraph{Learning, and teaching} A result presented three times, as a proof written for human readers, as a checked statement and as a
calculation, is a teaching tool: the student sees the same idea from three angles, and also sees where each one stops.
That is the bet of the textbook in Chapter~\ref{ch:c08}. It is also the bet of this book, at a simpler level.

\paragraph{Questions that touch the real world} The twenty contributions talk about real materials (crystals of charged
dust, semiconductors one atom thick, superfluid helium), about the history of the universe and about the evolution of
populations. None of them cures a disease or produces energy. But each shows that a concrete question can be linked to
a checked proof, without cheating about what remains to be done.

\begin{analogie}[A picture: the apprentice's notebook]
An apprentice who writes down every evening what he got right, what he got wrong and what he cannot yet do is not yet a
master. But his notebook is worth its weight in gold to whoever comes after him. This book is a notebook of that kind.
The picture has its limit: an apprentice's notebook is private, whereas this one is public, with everything that
obliges it to say.
\end{analogie}

\section{What remains out of reach}

We have to say honestly what this path does not reach, or not yet.

\paragraph{No human mathematician has reviewed these results} This is the most important limit. Everything at level~A
was checked by the kernel; nothing has been reviewed by a human specialist who would say: ``yes, this statement does
say what you claim''. That is why almost all the ledgers end with ``no''.

\paragraph{The dependence on openai/math} A large part of the book rests on a theorem produced by an AI from OpenAI,
checked by computer but not peer reviewed, and whose statement no person has audited. If that statement did not say
exactly what we think it says, Chapters~\ref{ch:c16}, \ref{ch:c03}, \ref{ch:c08}, \ref{ch:c09}, \ref{ch:c04} and~\ref{ch:c10} would
lose their physical significance, even though their logic would remain correct.

\paragraph{Scale} The judge checks what it is given, at the pace it can manage. Some certificates are too big; some
calculations are done at a single point only (in Chapter~\ref{ch:c01}, a single value of the matter density); some
modules have not been replayed yet.

\paragraph{Human reach} Are these results useful to researchers, to engineers, to teachers? We hope so, for the
textbook (Chapter~\ref{ch:c08}) and for the method. We do not know.

\section{What we do not know yet}

\begin{encadre}[Open questions at the end of the year]
\begin{itemize}
\item Do the level~A Lean statements say exactly what their glosses in plain language say? Only a human audit of the
statements can answer.
\item Is the upstream openai/math theorem correct in the sense mathematicians mean, and is its statement faithful to
the conjecture it claims to settle?
\item Is the triangular lattice the best for a real quantum crystal, and not only for the simplified model of
Chapters~\ref{ch:c12} and~\ref{ch:c11}?
\item Can the eight remaining stages of the Sandier--Serfaty certificate be brought into the kernel at a reasonable
cost?
\item Does the loop hold when you go faster, or when there are several of you? The October relapse says that it can
give way.
\item How can we organise, for results produced by machines, a human review that is neither a formality nor a
bottleneck?
\end{itemize}
\end{encadre}

\section{A science that makes sense}

I would like to end with what gave this year its meaning. It was not producing twenty contributions: it was learning,
the hard way, to tell apart what we know from what we believe. Today's machines write better than many of us; they do
not know, on their own, what is worth looking for, nor when to say ``I was wrong''. The collaboration that worked is
not the one in which the human lets the machine get on with it, nor the one in which he distrusts it so much that he
gets nothing out of it. It is the one in which each does what it does best: the machine proposes and checks, without
tiring; the human chooses, doubts, and answers for what is published.

If this way of working can be useful to others, beyond lattices and galaxies, then this book will have been useful.
That is a hypothesis. It is only level~C.
